\chapter{System Analysis and Design}
\label{ch:design}

This chapter explains the design of the ride-sharing application through its requirement analysis, high-level architecture, matching process, database design and class diagrams. Where a decision responds to a finding from the background literature (Chapter~\ref{ch:background}), the finding, the decision and the source are stated together.

%---------------------------------------------------------------
\section{Requirement Analysis}

The system requirements were analysed using Use Case modelling to identify the interactions between the user actors (Driver and Passenger) and the system. Only the major and complex use cases are detailed below, as they directly inform the identification of the system's core classes and components.

\begin{figure}[htbp]
  \centering
  \figbox{images/fig1-usecase.pdf}
  \caption{Use Case Diagram for the ride-sharing system}
  \label{fig:usecase}
\end{figure}

\subsection{Major Use Case Scenarios}

\paragraph{Use Case 1: Requesting a Ride and Matching (Complex)}
\begin{itemize}
  \item \textbf{Actor:} Passenger
  \item \textbf{Preconditions:} The passenger is authenticated and has submitted a ride request (RQ04).
  \item \textbf{Main Flow:}
  \begin{enumerate}
    \item The system fetches all active driver rides from the database.
    \item For each candidate ride, the system calculates the route overlap and detour distance using the Google Maps Routes API.
    \item The system applies the matching algorithm to filter out rides that exceed the maximum detour or time window thresholds.
    \item Compatible rides are scored and ranked based on proximity, time compatibility, and route overlap (RQ05).
    \item The system returns the ranked list of matches to the passenger.
  \end{enumerate}
  \item \textbf{Postconditions:} A list of \texttt{MatchResult} records is generated and displayed to the passenger.
\end{itemize}

\paragraph{Use Case 2: Posting a Ride (Major)}
\begin{itemize}
  \item \textbf{Actor:} Driver
  \item \textbf{Preconditions:} The driver is authenticated.
  \item \textbf{Main Flow:}
  \begin{enumerate}
    \item The driver inputs the origin, destination, departure time, and available seats.
    \item The system calls the Google Maps Routes API to calculate the route geometry, distance, and estimated duration.
    \item The system stores the ride offer and the calculated route data in the database (RQ03).
  \end{enumerate}
  \item \textbf{Postconditions:} A new \texttt{Ride} record is created and made visible to passengers.
\end{itemize}

\subsection{Identification of Classes from Use Cases}

The analysis of the above use cases directly informed the system's class structure. For example, the ``Requesting a Ride'' use case involves comparing a passenger's journey against multiple driver journeys, which identified the need for a dedicated \texttt{MatchingService} class to handle the computation, and a \texttt{MatchResult} class to store the output. Similarly, the ``Posting a Ride'' use case requires external API calls, identifying the need for a \texttt{RouteService} class to encapsulate the Google Maps API communication.

%---------------------------------------------------------------
\section{Overview of Design}

\subsection{System Architecture}

The high-level architecture is shown in Figure~\ref{fig:architecture}. The system has three parts: the Expo client application on the user's mobile device; the Supabase backend, which provides authentication, data storage and server-side logic; and the Google Maps Routes API (Application Programming Interface), an external service that supplies road-network route information.

\begin{figure}[htbp]
  \centering
  \figbox{images/fig2-architecture.pdf}
  \caption{High-level architecture of the ride-sharing system}
  \label{fig:architecture}
\end{figure}

The client and the backend communicate over REST (Representational State Transfer) for request--response calls, such as posting a ride, and over a realtime channel that pushes changes, such as the status of a request, to the client \citep{supabase2026}. Unlike the on-demand ride-hailing platforms reviewed in Section~\ref{sec:existing}, where a driver is assigned after a passenger requests a ride, the drivers in this system have already planned their journeys. Each journey is therefore stored in advance as a ride offer, and the backend compares complete journeys when a passenger submits a request (Section~\ref{sec:matching}). This requires the server-side matching and routing components described below.

\begin{itemize}
  \item \textbf{Finding:} On-demand ride-hailing platforms suffer from ``deadheading,'' where drivers travel empty to pick up a passenger, increasing costs and congestion \citep{furuhata2013}.
  \item \textbf{Decision:} Unlike on-demand systems, drivers in this application have pre-planned their journeys. Each journey is stored in advance as a ride offer, and the backend compares complete journeys only when a passenger submits a request.
  \item \textbf{Reference:} This design eliminates deadheading, aligning with the peer-to-peer cost-sharing model defined by \citet{furuhata2013}.
\end{itemize}

\subsection{Architectural Components}

\paragraph{Expo client.} The client uses React Native and Expo Router so that one code base serves both Android and iOS \citep{expo2026,meta2026}, which addresses RQ14. Its user interface (UI) screens (home, rides, profile and authentication) let users register and sign in, switch role, post, search and request rides, respond to requests, view status and history, and select locations on the interactive map (RQ01--RQ04, RQ06--RQ09, RQ11). The screens contain presentation logic only. Between the screens and the backend sits a thin service layer (Ride, Request and Profile services), which forwards each action to Supabase as a CRUD (Create, Read, Update, Delete) call. This separation supports the modularity requirement (NFR08), because a change in how data is stored does not require changes to the screens.

\paragraph{Supabase backend.} The backend is built on a PostgreSQL database and provides authentication, storage and server-side functions without a separate server application \citep{supabase2026}. Auth manages registration, sign-in and sessions; a JWT (JSON Web Token) issued at sign-in identifies the user in later requests, so protected functionality and data are restricted to the authenticated user and their role (RQ01, RQ13, NFR06). The database stores profiles, rides, requests and match results and supports ride histories (RQ08, RQ12).

\paragraph{Edge function.} The Edge function holds the matching service and the route service. Matching compares a passenger's journey with several candidate rides using route information from an external API (Section~\ref{sec:matching}; \citealp{google2026a}), so this logic was placed on the server and not in the client. This ensures that the same matching criteria apply on every device (RQ15, NFR08) and keeps the credentials for the routing service on the server, off the mobile device (NFR06).

\paragraph{Google Maps Routes API.} The Routes API returns the route distance, the estimated duration and an encoded polyline of the route geometry \citep{google2026a}. The distance between two coordinates does not represent a vehicle journey, whereas road-based values do (Section~\ref{sec:geo}), and route overlap indicates compatibility better than user locations alone \citep{hong2017}. Route information is therefore obtained from the API and not calculated from coordinates. Because each call adds overhead and the system depends on the service being available (NFR07), route calculations are limited to the journeys the matching process needs (Section~\ref{sec:geo}; \citealp{google2026b}).

\begin{itemize}
  \item \textbf{Finding:} Calculating road-based route overlap indicates compatibility significantly better than straight-line user locations \citep{hong2017}.
  \item \textbf{Decision:} The matching logic was placed on the server-side Edge function rather than the client, and it utilises the Google Maps Routes API to obtain accurate route geometry.
  \item \textbf{Reference:} This ensures consistent matching criteria across all devices (RQ15) and keeps external API credentials secure on the server (RQ13), as recommended by \citet{google2026a}.
\end{itemize}

\subsection{Matching Process}

The matching process is shown in Figure~\ref{fig:matching}. When a passenger submits a request (RQ04), the active driver rides are fetched as candidates. The route of each candidate is compared with the passenger's journey \citep{hong2017} and the ride is added to the match list if the overlap meets the threshold, or discarded if it does not. The comparison repeats until every candidate has been evaluated, and the matches are then ranked and returned (RQ05, RQ15). Unsuitable rides are therefore eliminated first and the remainder ranked, as described in Section~\ref{sec:matching}.

\begin{figure}[htbp]
  \centering
  \figbox{images/fig3-matching-flowchart.pdf}
  \caption{Flow chart of the ride-matching process}
  \label{fig:matching}
\end{figure}

In the class design (Section~\ref{sec:serviceclasses}), the overlap test is realised by \texttt{MatchingService}: a ride is accepted only when it has enough free seats and its pickup distance, detour and difference in departure time are within set limits. Accepted rides are scored and returned as \texttt{MatchResult} records.

%---------------------------------------------------------------
\section{Entity Relationship Diagram (ERD)}

The database layout is illustrated in Figure~\ref{fig:erd}. All tables required for the system are shown in the diagram, but only the major tables addressing core requirements are discussed below.

\begin{figure}[htbp]
  \centering
  \figbox{images/fig4-erd.pdf}
  \caption{Entity Relationship Diagram (ERD) of the database}
  \label{fig:erd}
\end{figure}

\begin{itemize}
  \item \textbf{Profile Table:} Represents a registered user (RQ02). It stores personal details and the user's current role (Driver or Passenger).
  \item \textbf{Ride Table:} Represents a ride offer posted by a driver (RQ03). It stores the origin, destination, departure time, and \texttt{seatsAvailable}. Crucially, it also stores the \texttt{routePolyline}, \texttt{routeDistanceM}, and \texttt{estDurationS} obtained from the Routes API when the ride is posted. Storing this data prevents redundant API calls during the matching process.
  \item \textbf{RideRequest Table:} Represents a passenger's request to travel (RQ04). It stores the requested origin, destination, time, and the current status of the request (e.g., Pending, Accepted, Declined).
  \item \textbf{MatchResult Table:} Stores the calculated compatibility scores between a specific \texttt{RideRequest} and candidate \texttt{Ride} records (RQ05).
\end{itemize}

The relationships are strictly defined: one \texttt{Profile} can post zero or many \texttt{Ride} records (1 to 0..*), and one \texttt{Profile} can submit zero or many \texttt{RideRequest} records (1 to 0..*). A \texttt{Ride} can receive zero or many \texttt{RideRequest} records, but the number of accepted requests is constrained by the \texttt{seatsAvailable} attribute.

%---------------------------------------------------------------
\section{Class Diagrams}

The classes are divided into entity classes, which represent stored data, and service classes, which hold the behaviour that creates, retrieves and processes it. The entity classes correspond to the database records in Figure~\ref{fig:erd}, and the service classes correspond to the thin service layer and the Edge function. The diagrams use UML (Unified Modelling Language) notation, in which \texttt{+} and \texttt{-} mark public and private members and multiplicities show how many objects of one class relate to one object of another. Parameter types are shown in the figures and omitted from the tables.

\subsection{Entity Classes}
\label{sec:entityclasses}

Figure~\ref{fig:entity} shows the four entity classes. Every entity has an \texttt{id} of type UUID (Universally Unique Identifier), locations are stored as latitude and longitude pairs (Section~\ref{sec:geo}), and a suffix gives the unit of a measured value (\texttt{M} for metres, \texttt{S} for seconds, \texttt{Min} for minutes).

\begin{figure}[htbp]
  \centering
  \figbox{images/fig3-entity-classes.pdf}
  \caption{Entity classes of the system}
  \label{fig:entity}
\end{figure}

\paragraph{Profile.} \texttt{Profile} represents a registered user (RQ02). One profile serves both drivers and passengers, with \texttt{role} recording the current role, so a user does not need to register twice. Credentials are not stored here because they are handled by Auth (RQ01).

\paragraph{Ride.} \texttt{Ride} is a ride offer posted by a driver for a planned journey (RQ03) and is the central entity for matching, search and history. It stores the origin, destination, departure time and \texttt{seatsAvailable}, together with the route (\texttt{routePolyline}, \texttt{routeDistanceM}, \texttt{estDurationS}) obtained from the Routes API when the ride is posted. Storing the route avoids requesting it again for every passenger request, which reduces external calls during matching (NFR01).

\paragraph{RideRequest.} \texttt{RideRequest} is a passenger's request to travel (RQ04). It stores the origin, destination and \texttt{requestedTime} used in matching, the \texttt{seatsNeeded} and the \texttt{status}. The \texttt{rideId} records the ride the passenger chooses and is set through \texttt{chooseRide()}. The operations \texttt{accept()}, \texttt{decline()}, \texttt{cancel()} and \texttt{complete()} represent the life cycle of a request (RQ06, RQ07).

\paragraph{MatchResult.} \texttt{MatchResult} records the outcome of comparing one request with one candidate ride (RQ05, RQ15). Its measurements, \texttt{pickupDistanceM}, \texttt{detourM} and \texttt{timeDiffMin}, correspond to compatibility factors identified in the literature (\citealp{stiglic2016}; Section~\ref{sec:factors}) and are combined into a score for ranking. It has no operations because it is a record; the behaviour that produces it belongs to \texttt{MatchingService}.

The operations of the entity classes are summarised in Table~\ref{tab:entityops}. The seat operations ensure that a ride cannot be allocated more passengers than the seats it offers (NFR03).

\begingroup\small
\begin{longtable}{@{}>{\raggedright\arraybackslash}p{0.15\textwidth}
                   >{\raggedright\arraybackslash}p{0.30\textwidth}
                   >{\raggedright\arraybackslash}p{0.47\textwidth}@{}}
\caption{Operations of the entity classes}\label{tab:entityops}\\
\toprule
\textbf{Class} & \textbf{Operation} & \textbf{Description} \\
\midrule
\endfirsthead
\toprule
\textbf{Class} & \textbf{Operation} & \textbf{Description} \\
\midrule
\endhead
\bottomrule
\endfoot
Profile & \texttt{switchRole()} & Changes the active role between driver and passenger. \\
Profile & \texttt{updateDetails(fullName)} & Updates the user's personal details. \\
Ride & \texttt{hasFreeSeat()} & Returns whether at least one seat remains; full rides are not offered to further passengers. \\
Ride & \texttt{takeSeats(n, passID)} & Allocates $n$ seats to a passenger and reduces \texttt{seatsAvailable}. \\
Ride & \texttt{releaseSeats(n)} & Returns $n$ seats, for example when a request is cancelled. \\
Ride & \texttt{cancel()}, \texttt{complete()} & Set the status of the ride to cancelled or completed. \\
RideRequest & \texttt{chooseRide(rideId)} & Links the request to the ride selected by the passenger. \\
RideRequest & \texttt{accept()}, \texttt{decline()} & Record the driver's decision on the request. \\
RideRequest & \texttt{cancel()}, \texttt{complete()} & Set the status of the request to cancelled or completed. \\
\end{longtable}
\endgroup

The multiplicities in Figure~\ref{fig:entity} are as follows. A profile posts zero or more rides as a driver and submits zero or more requests as a passenger (1 to 0..* in both cases), and each ride and request belongs to exactly one profile. A ride may receive zero or more requests (1 to 0..*), although the number that can be accepted is limited by \texttt{seatsAvailable}. A request produces zero or more match results, one for each compatible ride, and a ride may appear in zero or more match results (1 to 0..* in both cases); each match result refers to exactly one request and one ride. These associations are implemented through the identifier attributes \texttt{driverId}, \texttt{passengerId}, \texttt{rideId} and \texttt{requestId}, which act as references (foreign keys) between database records. To ensure all requirements are accounted for, these multiplicities are explicitly outlined to demonstrate that the system can handle multiple concurrent requests and matches without data loss, strictly adhering to the 1 to 0..* relationships defined above.

\subsection{Service Classes}
\label{sec:serviceclasses}

Figure~\ref{fig:service} shows the service classes. An arrow means that the class at its tail uses the class at its head, and colours link each service to its entity (orange for \texttt{Profile}, green for \texttt{Ride}, blue for \texttt{RideRequest}, pink for \texttt{MatchResult}; grey for infrastructure). Matching is a process and not stored data, so it is modelled as a service class, \texttt{MatchingService}, with \texttt{MatchResult} storing its output. Ride-matching combines several constraints and cannot be reduced to a single measure \citep{agatz2012,stiglic2016}, and this class keeps them in one place.

\begin{figure}[htbp]
  \centering
  \figbox{images/fig4-service-classes.pdf}
  \caption{Service classes of the system}
  \label{fig:service}
\end{figure}

\paragraph{SupabaseClient and ProfileService.} \texttt{SupabaseClient} is the single access point to the backend. Its operations \texttt{signUp()}, \texttt{signIn()} and \texttt{signOut()} handle authentication (RQ01, RQ13), \texttt{query()} reads records (RQ12), and \texttt{subscribe()} registers a callback for changes in a table so that status updates arrive in realtime (RQ07, RQ10). \texttt{ProfileService}, \texttt{RideService} and \texttt{RequestService} depend on it and not on the backend directly, so a change in backend access affects one class (NFR08). \texttt{RouteService} does not use \texttt{SupabaseClient} because it only communicates with the Google Maps Routes API. \texttt{ProfileService} retrieves and updates the signed-in user's profile through \texttt{SupabaseClient} (RQ02); its operations mirror those of \texttt{Profile} but locate the profile by \texttt{id} and save the change.

\paragraph{RideService and RequestService.} These classes provide the operations used to work with ride offers and requests, listed in Table~\ref{tab:serviceops}. \texttt{RideService} depends on \texttt{RouteService} to obtain a route when a ride is posted; because \texttt{RouteService} runs on the server, this dependency crosses from the client to the backend over the REST channel. \texttt{MatchingService} depends on \texttt{RequestService} because matching is performed on stored requests.

\begingroup\small
\begin{longtable}{@{}>{\raggedright\arraybackslash}p{0.19\textwidth}
                   >{\raggedright\arraybackslash}p{0.36\textwidth}
                   >{\raggedright\arraybackslash}p{0.37\textwidth}@{}}
\caption{Operations of \texttt{RideService} and \texttt{RequestService}}\label{tab:serviceops}\\
\toprule
\textbf{Class} & \textbf{Operation} & \textbf{Description} \\
\midrule
\endfirsthead
\toprule
\textbf{Class} & \textbf{Operation} & \textbf{Description} \\
\midrule
\endhead
\bottomrule
\endfoot
RideService & \texttt{postRide(driverId, journey, seats)} & Returns the new \texttt{Ride}. Creates a ride offer from the driver's journey and seats and stores its route (RQ03). \\
RideService & \texttt{cancelRide(rideId)}, \texttt{completeRide(rideId)} & Cancel or complete a ride (RQ07). \\
RideService & \texttt{searchRides(filters)} & Returns available rides matching criteria such as destination, departure time and seats (RQ11). \\
RideService & \texttt{getRideHistory(userId)} & Returns the user's previous rides (RQ08). \\
RequestService & \texttt{submitRequest(passengerId, journey)} & Returns the new \texttt{RideRequest}; matching is then performed for it (RQ04, RQ05). \\
RequestService & \texttt{chooseRide(requestId, rideId)} & Links the request to the ride the passenger selected. \\
RequestService & \texttt{acceptRequest(requestId)} & Records the driver's acceptance; the seats are then allocated with \texttt{takeSeats()} (RQ06). \\
RequestService & \texttt{declineRequest(requestId)} & Records the driver's refusal (RQ06). \\
RequestService & \texttt{cancelRequest(requestId)} & Cancels the request; allocated seats are returned with \texttt{releaseSeats()}. \\
RequestService & \texttt{getRequestsForRide(rideId)} & Returns the requests for a ride so the driver can respond (RQ06). \\
RequestService & \texttt{getRequestHistory(userId)} & Returns the user's previous requests (RQ08). \\
\end{longtable}
\endgroup

\paragraph{MatchingService.} \texttt{MatchingService} implements the process in Figure~\ref{fig:matching} (RQ05, RQ15). Its private attributes \texttt{maxPickupDistanceM}, \texttt{maxDetourM} and \texttt{timeWindowMin} are limits derived from the compatibility factors in Section~\ref{sec:factors}. Flexibility in departure time and detour influences the number of matches, but too much flexibility reduces a ride's attractiveness \citep{stiglic2016}, so the limits are adjustable attributes that can be tuned during testing (Section~\ref{sec:matching}). Only \texttt{findMatches()} is public; the other operations are private internal steps (Table~\ref{tab:matchingops}).

\begin{itemize}
  \item \textbf{Finding:} Flexibility in departure time and detour influences match rates, but too much flexibility reduces a ride's attractiveness \citep{stiglic2016}.
  \item \textbf{Decision:} \texttt{MatchingService} includes private attributes for \texttt{maxPickupDistance}, \texttt{maxDetour}, and \texttt{timeWindow} which act as strict filtering thresholds before scoring.
  \item \textbf{Reference:} This ensures the algorithm balances match quantity with driver convenience, as advised by \citet{stiglic2016}.
\end{itemize}

\begingroup\small
\begin{longtable}{@{}>{\raggedright\arraybackslash}p{0.42\textwidth}
                   >{\raggedright\arraybackslash}p{0.50\textwidth}@{}}
\caption{Operations of \texttt{MatchingService}}\label{tab:matchingops}\\
\toprule
\textbf{Operation} & \textbf{Description} \\
\midrule
\endfirsthead
\toprule
\textbf{Operation} & \textbf{Description} \\
\midrule
\endhead
\bottomrule
\endfoot
\texttt{+ findMatches(request)} & Returns a list of \texttt{MatchResult}. Evaluates each candidate ride, discards incompatible rides and returns the rest ranked by score. \\
\texttt{- isCompatible(ride, request)} & Returns true when enough seats are free for \texttt{seatsNeeded} and the pickup distance, detour and time difference are within the limits. \\
\texttt{- calcPickupDistance(ride, request)} & Returns the distance between the pick-up location and the route of the ride. \\
\texttt{- calcDetour(ride, request)} & Returns the additional distance the driver would travel to accommodate the passenger. \\
\texttt{- scoreMatch(ride, request)} & Combines the measurements into one score used for ranking. \\
\end{longtable}
\endgroup

\paragraph{RouteService.} \texttt{RouteService} encapsulates communication with the Google Maps Routes API and holds the route values returned: \texttt{distanceM}, \texttt{durationS} and \texttt{polyline} \citep{google2026a}. \texttt{RideService} uses it when a ride is posted and \texttt{MatchingService} when routes are compared. Isolating the external service in one class means that only this class changes if the routing provider changes (NFR08).

\subsection{Requirements Traceability}

Table~\ref{tab:traceability} shows how each requirement in the requirements catalogue (Appendix~\ref{app:catalogue}) is realised in the design.

\begingroup\small
\begin{longtable}{@{}>{\raggedright\arraybackslash}p{0.09\textwidth}
                   >{\raggedright\arraybackslash}p{0.3\textwidth}
                   >{\raggedright\arraybackslash}p{0.55\textwidth}@{}}
\caption{Traceability of requirements to components and classes}\label{tab:traceability}\\
\toprule
\textbf{ID} & \textbf{Requirement} & \textbf{Realised by} \\
\midrule
\endfirsthead
\toprule
\textbf{ID} & \textbf{Requirement} & \textbf{Realised by} \\
\midrule
\endhead
\bottomrule
\endfoot
RQ01 & Registration and authentication & Auth; \texttt{SupabaseClient} (\texttt{signUp}, \texttt{signIn}, \texttt{signOut}) \\
RQ02 & User roles and profiles & \texttt{Profile}; \texttt{ProfileService} \\
RQ03 & Driver ride posting & \texttt{Ride}; \texttt{RideService.postRide} \\
RQ04 & Passenger ride request & \texttt{RideRequest}; \texttt{RequestService.submitRequest} \\
RQ05 & Ride matching & \texttt{MatchingService}; \texttt{MatchResult}; \texttt{RouteService} \\
RQ06 & Ride request management & \texttt{RideRequest} (\texttt{accept}, \texttt{decline}); \texttt{RequestService} \\
RQ07 & Ride status management & \texttt{status} attributes; cancel and complete operations; \texttt{SupabaseClient.subscribe} \\
RQ08 & Ride history & \texttt{RideService.getRideHistory}; \texttt{RequestService.getRequestHistory} \\
RQ09 & Interactive map & UI screens; coordinate attributes; \texttt{RouteService.polyline} \\
RQ10 & Notifications & \texttt{SupabaseClient.subscribe} (realtime channel) \\
RQ11 & Ride search and filtering & \texttt{RideService.searchRides} \\
RQ12 & Data storage and retrieval & Database; \texttt{SupabaseClient.query}; all entity classes \\
RQ13 & Access control & Auth (sessions, JWT); \texttt{Profile.role} \\
RQ14 & Mobile application support & Expo client application (React Native) \\
RQ15 & Ride matching results & \texttt{MatchingService.findMatches}; \texttt{MatchResult} \\
RQ16 & Core task completion & UI screens (usability evaluated with the System Usability Scale) \\
\end{longtable}
\endgroup
